\subsection{测度的上包络与和}

命题 6. --- 设 \((\lambda_{\alpha})_{\alpha \in A}\) 是 T 上的一个递增定向的测度族，令 \(p = \sup_{\alpha} \lambda_{\alpha}^{\bullet}\) 。为使族 \((\lambda_{\alpha})\) 在 \(\mathcal{M}(T)\) 中有上界，必须且只需负担 \(p\) 局部有界。此时族 \((\lambda_{\alpha})\) 有上确界 \(\lambda\) （在 \(\mathcal{M}(T)\) 中），且 \(\lambda^{\bullet} = p\) 。对每个紧集 K，测度 \(\lambda_{K}\) 都是诸测度 \((\lambda_{\alpha})_{K}\) 在 \(\mathcal{M}(K)\) 中的上确界。

若族 \((\lambda_{\alpha})\) 在 \(\mathcal{M}(\mathrm{T})\) 中有上界，则 p 显然局部有界。反之，设 p 局部有界，我们来证明它满足 No.~2 的 Prop. 2 b) 的条件 1) 与 2)。对于 2)，这由下列等式得出：

\[
p (f) = \sup _ {\alpha} \lambda_ {\alpha} ^ {\bullet} (f) = \sup _ {\alpha} \sup _ {K} \lambda_ {\alpha} ^ {\bullet} (f \varphi_ {K}) = \sup _ {K} \sup _ {\alpha} \lambda_ {\alpha} ^ {\bullet} (f \varphi_ {K}) = \sup _ {K} p (f \varphi_ {K}).
\]

另一方面，设 K 是一个紧集；负担 \(p_{K}\) 等于诸负担 \((\lambda_{\alpha}^{\bullet})_{K}\) 的上包络，且由于 p 局部有界，它是有界的。于是诸测度 \((\lambda_{\alpha})_{K}\) 有上确界 \(\lambda_{K}\) （在 \(\mathcal{M}(K)\) 中），且 \(\lambda_{K}^{\bullet}=p_{K}\) (Ch. V, §1, No.~4, Prop. 11)。这样 No.~2 的 Prop. 2 b) 的条件 1) 得以满足，因此存在 T 上的测度 \(\lambda\) 使 \(\lambda^{\bullet}=p\) ；显然 \(\lambda\) 是诸测度 \(\lambda_{\alpha}\) 的上确界。

定义 9. --- 设 \((\mu_i)_{i \in I}\) 是 T 上的一个测度族。设 A 是 I 的有限子集全体；对每个 \(\alpha \in A\) 令 \(\lambda_\alpha = \sum_{i \in \alpha} \mu_i\) 。若族 \((\lambda_\alpha)\) 有上确界 \(\mu\) （在 \(\mathcal{M}(T)\) 中），则称族 \((\mu_i)\) 可和，\(\mu\) 称为族 \((\mu_i)\) 的和，并记 \(\mu = \sum_{i \in I} \mu_i\) 。

这一定义推广了 Ch. V, §2, No.~1 中的定义。

命题 7. --- 为使族 \((\mu_i)_{i \in I}\) 可和且其和为 \(\mu\) ，必须且只需负担 \(p = \sum_{i \in I} \mu_i^\bullet\) 局部有界，此时 \(p = \mu^\bullet\) 。于是对 T 的每个紧子集 K，族 \(\left((\mu_i)_K\right)_{i \in I}\) 在 \(\mathcal{M}(K)\) 中可和，且 \(\mu_K = \sum_{i \in I} (\mu_i)_K\) 。

沿用 Def. 9 的记号，有 \(\lambda_{\alpha}^{\bullet} = \sum_{i\in \alpha}\mu_{i}^{\bullet}\) （对 I 的每个有限子集 \(\alpha\)）(No.~2, 注 1)。于是所述结论是 Prop. 6 的直接推论。

关系式 \(\mu_{\mathrm{K}} = \sum_{i\in \mathrm{I}}(\mu_i)_{\mathrm{K}}\) 连同 Ch. V, §2, No.~2 的 Prop. 2 给出下列结果：

命题 8. --- 设 \(\mu\) 是 T 上一可和测度族 \((\mu_i)_{i \in I}\) 的和。为使 T 到拓扑空间 F（Hausdorff 与否均可）中的映射 \(f\) 是 \(\mu\) -可测的，必须且只需 \(f\) 是 \(\mu_i\) -可测的（对每个 \(i \in I\)）。

